% Book pp. 62-63 (PDF pp. 63-64)
\section{Arithmetic and Geometric Means of Sequences}

The mean of any two numbers, for example, a and b, is obtained by finding a half
of the sum of the numbers thus: \textit{Mean}, $(\mu) = \frac{1}{2}(a + b)$. The difference
between a and $\mu$ is equal to the difference, or distance, between $\mu$ and b i.e.,
$\mu - a = b - \mu$.

Now, $a$, $\mu$, $b$ form an arithmetic sequence since there is a common difference and
that for any arithmetic sequence with a and b as terms, some means namely,
$\mu_1, \mu_2, \mu_3, \ldots \mu_k$ can be found between a and b such that a,
$\mu_1, \mu_2, \mu_3, \ldots \mu_k$, b, remains an arithmetic sequence as
$\mu_1 - a = \mu_2 - \mu_1 = \ldots \mu_k - b$ and the means are equally spaced.

\begin{example}{Example 2.11}
\begin{enumerate}
  \item Find the arithmetic mean sequence of students' scores in mathematics test:
  80, 75, 90, 85, 95.
  \item A baker wants to increase the amount of sugar in a recipe from 200g to
  500g in four equal steps. What arithmetic sequence represents the amount
  of sugar (in grammes) added and what are the intermediate values?
\end{enumerate}

\solution
\begin{enumerate}
  \item Arithmetic mean $= \dfrac{80 + 75 + 90 + 85 + 95}{5} = \dfrac{425}{5} = 85$
  \item The baker wants to increase the sugar from 200g to 500g
  \begin{align*}
    500 - 200 &= 300g \\
    d &= \frac{300}{4 - 1} \\
    &= \frac{300}{3} \\
    &= 100g
  \end{align*}
  Arithmetic mean
  \begin{align*}
    200 + 100 &= 300g \\
    300 + 100 &= 400g \\
    400 + 100 &= 500g
  \end{align*}
  The arithmetic sequence: 200, 300, 400, 500.

  The intermediate sequence 300 \textit{and} 400.
\end{enumerate}
\end{example}

\subsection{Geometric mean}

The geometric mean of a sequence is the measure of the central tendency of the
sequence, similar to the arithmetic mean. However, it is calculated differently.

For a sequence of $n$ numbers $(a_1 \times a_2 \times a_3 \ldots a_n)^{\frac{1}{n}}$

\begin{example}{Example 2.12}
\begin{enumerate}
  \item Find the geometric mean of the sequence 2, 4, 8, 16
  \item Find the geometric mean of the sequence 10, 20, 40.
\end{enumerate}

\solution
\begin{enumerate}
  \item 2, 4, 8, 16
  \[
    (a_1 \times a_2 \times a_3 \ldots a_n)^{\frac{1}{n}}
  \]
  \[
    (2 \times 4 \times 8 \times 16)^{\frac{1}{4}} = \sqrt[4]{1024} = 5.65685\ldots \approx 5.7
  \]
  \item $(10 \times 20 \times 40)^{\frac{1}{3}} = \sqrt[3]{8000} = 20$
\end{enumerate}
\end{example}
